%% Appendix -- the skills, verbatim, so students can copy them into skills/.
%% One frame per skill, grouped Plan (Unit 3) then Build (Unit 4). Each frame
%% reproduces skills/<name>.md line for line in a tdcodet box (the same dark
%% mono box the GLOSSARY.md frame uses), so what is on the slide IS the file.
%% Source of truth: setup/sandbox/skills/*.md. If a skill file changes, change
%% the matching frame here. No preamble; \input by the master after unit4.

% ======================================================================
\section[Skills]{The Skills, to Copy}
% ======================================================================

% -------------------------------------------------- how to use this appendix
\begin{frame}{Copy these into \texttt{skills/}}
  \begin{withmargin}
    \begin{tdmain}
      Every skill is just a markdown file the agent reads when you name it.
      Nothing here is magic. To get the whole toolkit, create one file per skill
      under \texttt{skills/} and paste the text on the next slides into it,
      keeping the \texttt{name:} and \texttt{description:} header as shown.
      \vskip0.8em
      \begin{tdblock}
        {\ttfamily\small
        skills/grill-me.md\\
        skills/domain-modeling.md\\
        skills/to-prd.md \;\ldots{} one file each}
      \end{tdblock}
    \end{tdmain}
    \begin{tdmargin}
      \alertbf{Plan} skills come first (Unit 3), then \alertbf{Build} skills
      (Unit 4).
      \par\medskip
      The filename must match the \texttt{name:} line; that is the word you
      invoke it by.
    \end{tdmargin}
  \end{withmargin}
\end{frame}

% ============================== PLAN SKILLS (Unit 3) ===================

% -------------------------------------------------------------- grill-me
% Longest skill: shrink to \tiny and let the box wrap each rule naturally so
% every line survives above the footer.
\begin{frame}{Plan skill: \texttt{grill-me}}
  \begin{tdcodet}{skills/grill-me.md}
    \tiny
    name: grill-me\\
    description: Interview the user relentlessly about a plan until you reach a shared understanding.\\[3pt]
    Interview me relentlessly about every aspect of this plan until we reach a shared understanding. The goal is not a document, it is alignment: a design concept we both hold.\\[3pt]
    Rules:\\[2pt]
    1. First, explore the codebase enough to ask informed questions. If a question can be answered by reading the code, read it instead of asking.\\
    2. Ask questions one at a time. Wait for my answer before the next one.\\
    3. For every question, propose your own recommended answer and say why.\\
    4. Walk down each branch of the decision tree, resolving dependencies one by one. Surface the awkward questions a non-expert would forget (edge cases, retroactive data, failure behaviour, out-of-scope boundaries).\\
    5. Do not write a plan or any code yet. Keep going until you have no open questions, then say so.\\[3pt]
    Stop when alignment is reached, not before.
  \end{tdcodet}
\end{frame}

% -------------------------------------------------------- domain-modeling
\begin{frame}{Plan skill: \texttt{domain-modeling}}
  \begin{tdcodet}{skills/domain-modeling.md}
    name: domain-modeling\\
    description: Build the shared vocabulary for the game and record it in GLOSSARY.md.\\[3pt]
    Pin down the words we use for the game so the code and tests all say the same thing.\\[3pt]
    1. List. Collect the nouns and verbs from the chat and the code: obstacle, hitbox, duck, ramp, spawn rate, run length.\\
    2. Define. Give each term one plain line: what it means and where it lives (createGame, nextState, collides).\\
    3. Record. Write these lines into GLOSSARY.md, one term per line.\\
    4. Fix. Find synonyms (two words, one thing) and clashes (one word, two things), then pick one winner and drop the rest.\\[3pt]
    One word, one meaning, everywhere.
  \end{tdcodet}
\end{frame}

% ---------------------------------------------------------------- to-prd
\begin{frame}{Plan skill: \texttt{to-prd}}
  \begin{tdcodet}{skills/to-prd.md}
    name: to-prd\\
    description: Turn an aligned conversation into a product requirements document.\\[3pt]
    Summarise the design concept we reached into a Product Requirements Document. Write it to issues/PRD.md.\\[3pt]
    Structure:\\[2pt]
    - Problem --- the problem the user faces, in one short paragraph.\\
    - Solution --- the approach, in one short paragraph.\\
    - User stories --- a short list of "as a <role>, I want <goal>, so that <reason>".\\
    - Implementation decisions --- the concrete choices we agreed on.\\
    - Testing decisions --- how each story will be verified.\\
    - Out of scope --- what we explicitly decided NOT to do, and why. This is the definition of done.\\[3pt]
    Do not write implementation code. The PRD records the destination, not the route.
  \end{tdcodet}
\end{frame}

% -------------------------------------------------------------- to-issues
\begin{frame}{Plan skill: \texttt{to-issues}}
  \begin{tdcodet}{skills/to-issues.md}
    \tiny
    name: to-issues\\
    description: Slice a PRD into independently grabbable vertical-slice issues.\\[3pt]
    Break the PRD in issues/PRD.md into independently grabbable issues, written as local markdown files in issues/ (one file per issue, e.g. issues/01-*.md).\\[3pt]
    Rules:\\[2pt]
    1. Slice vertically, not horizontally. A good slice cuts through every layer (data, service, dashboard) and produces something visible and testable. A bad slice is "all the schema" or "all the dashboard".\\
    2. The first slice must be a tracer bullet: the thinnest end-to-end path that proves the whole flow works.\\
    3. Record blocking relationships between issues (which must be done before which), so independent issues could be worked in parallel.\\
    4. Each issue states: title, the user-facing outcome, and how it will be tested.\\[3pt]
    If your first proposed slice is horizontal, reject it and re-slice before writing the files.
  \end{tdcodet}
\end{frame}

% -------------------------------------------------------------- prototype
\begin{frame}{Plan skill: \texttt{prototype}}
  \begin{tdcodet}{skills/prototype.md}
    name: prototype\\
    description: Build a throwaway prototype to answer one design question, then delete it.\\[3pt]
    Answer one design question with a throwaway prototype, then delete it.\\[3pt]
    1. Ask. Write the one design question in a single sentence (e.g. tick-based core or event-based?). If you cannot fit it on one line, stop and narrow it.\\
    2. Scratch. Build the smallest scratch script that answers only that question. Skip edge cases, node --test, and GLOSSARY.md vocabulary.\\
    3. Record. Write the answer as one line in the matching issues/ file, or in GLOSSARY.md if it changes the shared model.\\
    4. Delete. Remove the scratch files so no throwaway code survives.\\[3pt]
    The answer is the deliverable, not the code.
  \end{tdcodet}
\end{frame}

% ============================== BUILD SKILLS (Unit 4) ==================

% ----------------------------------------------------------------- tdd
\begin{frame}{Build skill: \texttt{tdd}}
  \begin{tdcodet}{skills/tdd.md}
    name: tdd\\
    description: Implement a change with the red, green, refactor loop.\\[3pt]
    Implement the task using test-driven development.\\[3pt]
    1. Red. Write or locate a single failing test that specifies the next small\\
    \phantom{1. }behaviour. Run node --test and confirm it fails for the right reason.\\
    2. Green. Write the minimum code to make that test pass. Run node --test.\\
    3. Refactor. With the bar green, tidy the code without changing behaviour.\\
    \phantom{3. }Run node --test again.\\[3pt]
    Make one small change at a time and run the tests after each. Do not edit the\\
    test to fit broken code; fix the code. When done, show the diff for review.
  \end{tdcodet}
\end{frame}

% -------------------------------------------------------- codebase-design
\begin{frame}{Build skill: \texttt{codebase-design}}
  \begin{tdcodet}{skills/codebase-design.md}
    name: codebase-design\\
    description: Design a deep module: a lot of behaviour behind a small, simple interface.\\[3pt]
    Design a deep module: hide a lot of behaviour behind a small, simple interface.\\[3pt]
    1. Name one job. Write one sentence for what this module does. One job only.\\
    2. Check GLOSSARY.md. Reuse the shared names from GLOSSARY.md so the interface fits the domain.\\
    3. Draft the small interface. Write only the signatures that hide the work, like createGame, nextState, collides. Nothing extra.\\
    4. Prove it is testable. Confirm each signature runs under node --test with no canvas and no DOM. If not, shrink it.\\
    5. Implement behind it. Now write the code inside those functions. Keep the interface fixed.\\[3pt]
    A deep module hides a lot behind a little.
  \end{tdcodet}
\end{frame}

% ------------------------------------------------------------- code-review
\begin{frame}{Build skill: \texttt{code-review}}
  \begin{tdcodet}{skills/code-review.md}
    name: code-review\\
    description: Review a diff twice, once for standards and once for spec, in a fresh context.\\[3pt]
    Read the diff fresh and report what is wrong.\\[3pt]
    1. Clear. Drop all prior context first, so you read the diff with no memory of writing it.\\
    2. Standards. Check the diff against the rules in GLOSSARY.md and AGENTS.md. The core in game.js must stay free of canvas and DOM.\\
    3. Spec. Check the diff does what the issue acceptance line or the PRD asked, no more and no less.\\
    4. Report. List every problem you found in one message.\\[3pt]
    Report the problems, do not rewrite the code.
  \end{tdcodet}
\end{frame}

% --------------------------------------------------------- diagnosing-bugs
\begin{frame}{Build skill: \texttt{diagnosing-bugs}}
  \begin{tdcodet}{skills/diagnosing-bugs.md}
    name: diagnosing-bugs\\
    description: Find and fix the root cause of a bug, not the symptom.\\[3pt]
    Trace a bug to its root cause in the pure core (createGame, nextState, collides) and fix that, not the surface symptom.\\[3pt]
    1. Reproduce. Write a failing test that triggers the bug. Run node --test and watch it fail.\\
    2. Minimise. Cut the test to the smallest input that still fails. Name what you see using GLOSSARY.md.\\
    3. Find the cause. State the likely cause in one sentence. Add a log to confirm it before you change anything.\\
    4. Fix. Change the cause, remove the log, and run node --test until green.\\[3pt]
    A patch that hides the cause creates two new bugs.
  \end{tdcodet}
\end{frame}

% ------------------------------------------- improve-codebase-architecture
\begin{frame}{Build skill: \texttt{improve-codebase-architecture}}
  \begin{tdcodet}{skills/improve-codebase-architecture.md}
    name: improve-codebase-architecture\\
    description: Find a tangled cluster and propose deepening it into one testable module.\\[3pt]
    Find one tangled cluster and propose deepening it into a single testable module.\\[3pt]
    1. Read. Open GLOSSARY.md for the shared vocabulary before you look at code.\\
    2. Scan. Find mixed concerns, like physics tangled into the render loop in index.html.\\
    3. Cut. Pick one clean boundary to pull into the pure core (createGame, nextState, collides).\\
    4. Propose. Name what moves behind game.js and its interface. Write the proposal only, do not edit code.\\[3pt]
    A boundary you can wrap a node --test around is the goal.
  \end{tdcodet}
\end{frame}
